\section{Related work}\label{sec:related}

\paragraph{Indian agricultural price forecasting.}
Most work comes from statistical and machine-learning groups modelling a few series in depth.
\begin{itemize}
  \item Paul et al.\ compared ARIMA, GRNN, SVR, random forests and gradient boosting on daily AGMARKNET brinjal prices
from 17 Odisha markets, with Diebold--Mariano tests and a model confidence set \cite{paul2022brinjal}.
  \item A later study of monthly wholesale prices for 19 commodities in 143 markets found deep recurrent models ahead
of ML and stochastic models. Unusually, it included a random-walk-with-drift baseline \cite{paul2025deep}.
  \item Wavelet hybrids, time-delay networks and ANN/linear hybrids have been proposed for oilseeds, pulses and other commodities
\cite{paul2021wavelet,jha2012tdnn,jha2013ann,anjoy2017pulses}.
  \item ARMA--GARCH models have been fitted to a single onion market \cite{kumar2021onion}, and LSTM variants to potato and other prices
\cite{praveenkumar2024potato,jaiswal2021lstm,manogna2021spot}.
  \item The closest study to ours in data coverage is Manogna et al.\ \cite{manogna2025deep}: 23 commodities, 165 markets, daily AGMARKNET data, 2010--2024.
    It uses an 80/20 split with sliding windows, reports RMSE and MAPE, and explicitly decides against MASE.
    Its abstract and body give different MAPE values for the same models.
  \item National CPI sub-indices for tomato, onion and potato have been forecast with autoregressive models \cite{vani2024top}.
  \item Reviews summarise the shift from statistical to intelligent and combined models \cite{sun2023review}.
\end{itemize}

\paragraph{Big-data systems.}
Spark \cite{zaharia2016spark} and MLlib \cite{meng2016mllib} are general platforms.
The nearest agricultural precedent we found uses Spark to clean multi-city prices in Pakistan and then models a single series \cite{rana2024bigdata}.
We found no system that trains, evaluates, calibrates and serves forecasts for thousands of mandi series within one distributed pipeline.

\paragraph{Volatility and anomaly detection.}
The IFPRI early-warning system flags excessive price variability in international futures with nonparametric thresholds on returns \cite{geng2021earlywarning}.
LSTM--GARCH hybrids have modelled volatility in Indian prices \cite{manogna2025volatility}.
We found no mandi-level system that combines point forecasts, intervals and anomaly flags.

\paragraph{Evaluation methodology and uncertainty.}
Our protocol draws on the following:
\begin{itemize}
  \item scaled errors \cite{hyndman2006another} and the Diebold--Mariano test \cite{diebold1995comparing};
  \item the M4 and M5 competitions, which show that simple methods are strong \cite{makridakis2020m4,makridakis2022m5};
  \item work on when cross-validation is valid for time series and on how to estimate performance \cite{bergmeir2018note,cerqueira2020evaluating};
  \item a catalogue of evaluation pitfalls \cite{hewamalage2022pitfalls}.
\end{itemize}
Conformal methods for time series include EnbPI \cite{xu2023conformal}, ensemble conformalised quantile regression \cite{jensen2024encqr} and adaptive conformal inference under shift \cite{gibbs2021adaptive}.
We use the simplest member of this family, split conformal with Mondrian bins, and \emph{measure} its coverage.

\paragraph{Market information and farmer welfare.}
\begin{itemize}
  \item Mobile phones reduced price dispersion among Kerala fishermen \cite{jensen2007digital}.
  \item Internet kiosks raised soy prices in Madhya Pradesh \cite{goyal2010information}.
  \item Randomised SMS price information in Maharashtra had small effects \cite{fafchamps2012sms}.
  \item Wholesale-price information for West Bengal potato farmers did not change average farm-gate prices \cite{mitra2018asymmetric}.
\end{itemize}
These results motivate our framing of decision cards as cautious market outlooks.

\paragraph{Gap.}
To our knowledge, none of the reviewed Indian studies jointly does all three of the following:
(i)~evaluates daily mandi-level modal prices across thousands of market series;
(ii)~uses strictly chronological, purged splits with naive and seasonal baselines and scale-free errors at several horizons;
(iii)~reports distribution-free intervals whose coverage is measured out of sample.
AG-DATA is designed to test whether doing all three changes the conclusions drawn from smaller studies.
